\section{Example Story}

%Keywords covered: 
%TODO Vulnerable
%TODO Blaze
%TODO Shield X
%TODO Lifepay X
%TODO Advantage X
%TODO Stun X
%TODO Poison X
%TODO Heal X
%TODO Powerful
%TODO Flurry

%Static effects covered
%TODO Ranged, TODO Selftargeting, TODO ally boons, TODO no crossmapping 
%TODO AOE
%TODO Multi Targeting
%TODO Elements
%TODO Vulnerable

%TODO passives, shaping, start of encounter, beginnign of turn

This section details an example round for how three players could go through the first story.
All keywords, boons \& banes will be covered, as well as all enemy behaviour and all types of passive.
This should suffice to prepare you for most situations.


\subsection{Setup}
Our four players are:
\begin{itemize}
	\item Alexander
	\item Brigitte
	\item Christian
	\item Dorothy
\end{itemize}

They start as all games must, with the setup.


\begin{figure}[ht]
	\begin{minipage}[t][6cm][t]{0.5\textwidth}
			\fbox{\parbox[4cm]{\textwidth}{Picture:  Human Koloss 1 Knight 1}}
	\end{minipage}
	\begin{minipage}[t][6cm][t]{0.5\textwidth}
		
		Alexander feels heroic today, and chooses to play a Human, Knight (Rank 1) and Koloss (Rank 1). His shuffles together the 16 action cards from his Ancestry and two classes. He takes the red and blue attack dice from the human, and sets up his Tableau like in the picture. The passive from Human changes his action deck. It contains the most relibly powerful card in the game: "Smite". The passive from Knight will allow him to see more cards at the start of the game, and will help find "Smite" and other good cards. His Koloss passive allows him to mitigate damage if it were to come in bursts
		
	\end{minipage}	
	
\end{figure}

%% Brigitte
\begin{figure}[ht]
	\begin{minipage}[t][6cm][t]{0.5\textwidth}
		Brigitte likes to play a slower controlled game, and chooses to play Fae, Witch (Rank 1) and Assassin (Rank1). She shuffles together the 15 action cards from her Ancestry and two classes. She takes the red and blue attack dice from Fae, and sets up her Tableau like in the picture. The passive from Fae allows her to always interact with her neighbours Alexanderand Christian in case they need a helping hand. The passive from Assassin lets her hold 5 cards in hand instead of 4, and also slims down her deck to 15 cards. This gives her more options in every turn. Her Witch passive lets her heal significantly when her enemies die, which will help drag the fight out when she assumes a winning position.
		
	\end{minipage}	
	\begin{minipage}[t][6cm][t]{0.5\textwidth}
		\fbox{\parbox[4cm]{\textwidth}{Picture:  Fae Assassin 1 Witch 1}}
	\end{minipage}
	
\end{figure}


%% Christian
\begin{figure}[ht]
	\begin{minipage}[t][5cm][t]{0.5\textwidth}
		\fbox{\parbox[4cm]{\textwidth}{Picture:  Halfelf Ranger 1 Bard 1}}
	\end{minipage}
	\begin{minipage}[t][9cm][t]{0.5\textwidth}
		Christian wants as many dice as he can have, and chooses to play Halfelf, Ranger (Rank1) and Bard (Rank 1). He shuffles together the 16 action cards from her Ancestry and two classes. He takes the two red and two blue attack dice from the Halfelf and one silver one from Ranger. He sets up his Tableau like in the picture. The passives from Halfelf and Ranger both grant him extra attack dice. His Bard passive adds two Wood to the wheel of elements at the start of each encounter and will give a head start on element production.
		
	\end{minipage}	
\end{figure}

	%% Dorothy
	\begin{figure}[ht]
		\begin{minipage}[t][9cm][t]{0.5\textwidth}
			Dorothy is here to lay down a beating,  and chooses to play Wyrmkin, Berserker (Rank1) and Druid (Rank 1). She shuffles together the 16 action cards from her Ancestry and two classes. She takes the red and blue attack dice from Wyrmkin, and sets up her Tableau like in the picture. The passive from Wyrmkin allows her sacrifice her life to deal more damage. The passive from Berserker lets her use shied to deal more damage instead of shielding. Her Druid passive heals her every turn to  offset the massive amount of life she can spend with her wyrmkin passive.
			
		\end{minipage}	
		\begin{minipage}[t][9cm][t]{0.5\textwidth}
			\fbox{\parbox[4cm]{\textwidth}{Picture:  Fae Assassin 1 Witch 1}}
		\end{minipage}
	\end{figure}


Since this is their first time playing, They decide to play the first encounter "Bandit Tracking"
Alexander reads out the story, and they procedd with the first encounter. 

\subsection{Encounter 1}

The players define Alexaner to be player 1, Brigitte is player 2, Christian is player 3 Dorothy is player 4. According to the encounter setup, thisordering does not really matter, since everyone gets a Bandit Skirmisher and two Bandit Raiders. 

Each player sets their starting life
Alexander sets his life to 

They roll the monster die, and it comes up on a 1.

Alexander draws four cards for the start of the turn


\subsection{Event \& Level up}

\subsection{Encounter 2}

\subsection{Event \& Level up}

\subsection{Encounter 3}